\documentclass[10pt,a4paper,onecolumn,article]{memoir}
\counterwithout{section}{chapter}
\DoubleSpacing

\usepackage[T1]{fontenc}
\usepackage{palatino}
\usepackage{graphicx}
  \graphicspath{{./figs/}}

%% remove colorlinks when preparing for print
\usepackage[final,hyperindex,hyperfootnotes,bookmarksnumbered,colorlinks]{hyperref}

\usepackage{eqparbox}
  %% keep size of first column of sequence align table on memory to calculate
  %% the size limit for the second column
  \newsavebox{\firstentrybox}
  \newcolumntype{F}{%
    >{\begin{lrbox}{\firstentrybox}}%
      l%
    <{\end{lrbox}%
    \eqmakebox[firstentry][l]{\unhcopy\firstentrybox}}}

\usepackage{color}    % to color background of generated values values
\usepackage{seqsplit}
\usepackage{url}
\usepackage[textsize=footnotesize]{todonotes}
  %% new command for box about missing references
  \newcommand{\addref}[1][]{\todo[color=red!40,size=\tiny]{Add reference: #1}}

\usepackage{natbib}

%% to keep a pair of related table and figures, we set a new float and use it
%%conjunction with the captionof command from capt-of
\usepackage{capt-of}
\newfloat{TableAndFigure}{X}{capname}

\input{results/variables-align_results}
\newcommand{\CodingGenesInHistOne}{50}
\newcommand{\PseudoGenesInHistOne}{9}
\newcommand{\TotalGenesInHistOne}{59}

\newcommand{\CodingGenesInHistTwo}{11}
\newcommand{\PseudoGenesInHistTwo}{7}
\newcommand{\TotalGenesInHistTwo}{XXXXX}

\newcommand{\CodingGenesInHistThree}{3}
\newcommand{\PseudoGenesInHistThree}{XXXX}
\newcommand{\TotalGenesInHistThree}{XXXXX}

\newcommand{\CodingGenesInHistFour}{1}
\newcommand{\PseudoGenesInHistFour}{XXXX}
\newcommand{\TotalGenesInHistFour}{XXXXX}

\newcommand{\HistOneSpan}{2.3\,Mbp}
\newcommand{\HistTwoSpan}{105\,kbp}
\newcommand{\HistThreeSpan}{XXXXbp}
\newcommand{\HistFourSpan}{XXXXbp}

\input{results/variables-protein_stats}
\input{results/variables-utr}
% 12 coding H2A, 10 coding H3, 6 coding H1
% 15 coding H2B, 12 coding H4
\newcommand{\MarzluffCodingGenesInHistOne}{50}
\newcommand{\MarzluffPseudoGenesInHistOne}{9}
\newcommand{\MarzluffTotalGenesInHistOne}{59}

% 3 coding H2A, 2 pseudo H3, 1 coding H3
% 4 pseudo H2B, 1 coding H2B, 1 coding H4
\newcommand{\MarzluffCodingGenesInHistTwo}{11}
\newcommand{\MarzluffPseudoGenesInHistTwo}{7}
\newcommand{\MarzluffTotalGenesInHistTwo}{XXXXX}

% 1 coding H2A, 1 coding H3
% 1 pseudo H2B, 1 coding H2B
\newcommand{\MarzluffCodingGenesInHistThree}{3}
\newcommand{\MarzluffPseudoGenesInHistThree}{XXXX}
\newcommand{\MarzluffTotalGenesInHistThree}{XXXXX}

% 1 coding H2A
\newcommand{\MarzluffCodingGenesInHistFour}{1}
\newcommand{\MarzluffPseudoGenesInHistFour}{0}
\newcommand{\MarzluffTotalGenesInHistFour}{0}

\newcommand{\MarzluffHistOneSpan}{2.3\,Mbp}
\newcommand{\MarzluffHistTwoSpan}{105\,kbp}
\newcommand{\MarzluffHistThreeSpan}{XXXXbp}
\newcommand{\MarzluffHistFourSpan}{XXXXbp}


%% We define a style here which will be used around all the values found
%% by the scripts. The idea is to place them in a grey box so that people
%% reading the paper can really get the idea that the values in the text
%% are automatically generated. The problem is that the variables are not
%% simple numbers or text, and we want to make operations with them here
%% in LaTeX. So we get an handy trick, and redefine FPeval.
%% See http://tex.stackexchange.com/questions/159155 for details
\newcommand{\ScriptValue}[1]{\setlength{\fboxsep}{2pt}\colorbox[gray]{0.8}{#1}}
\let\OldFPeval\FPeval
\def\FPeval{\renewcommand{\ScriptValue}[1]{##1}\OldFPeval}

\author{David Miguel Susano Pinto and Andrew Flaus}
\title{The canonical core histone catalogue}

\begin{document}
  %% Text flow: first introduction on what are histones and nucleosomes. Some history
  %% on them, and explanation on the different types. Then we list their organization,
  %% following the same pattern of gene expression: first describe genes and their
  %% organization, then their transcripts, and finally the proteins. One everything is
  %% presented, discuss on how the current division of histones is inappropriate and
  %% leads to misconceptions about histone variance. Finalize with comment on reproducible
  %% research and the code availability to rebuild all tables and analysis of this paper,
  %% and how we have used it for other organisms.

  \maketitle

  \begin{abstract}
	Core histone proteins are encoded by a large family of genes distributed
	in \todo{total histone clusters} clusters across the human genome.
	In total there are \TotalGenes{} canonical human core histone genes, including
	\HTwoACodingGenes{} for H2A, \HTwoBCodingGenes{} for H2B,
	\HThreeCodingGenes{} for H3 and \HFourCodingGenes{} for H4, 
	along with at least \TotalPseudoGenes{} total pseudogenes.
	The histone genes display coding variation that gives rise to
	\HTwoAUniqueProteins{} H2A, \HTwoBUniqueProteins{} H2B, 
	\HThreeUniqueProteins{} H3 and \HFourUniqueProteins{} unique protein isoforms.
	This represents a surprising and seldom recognised variation
	of up to \todo{highest identity}\% amino acid identity. 
	The gene sequence and protein isoform diversity exists 
	in parallel to the well-known specialisation of variant histones 
	and exceeds commonly used designations of variation such as H2A.1 and H3.1. 
	As a first step to clarifying whether specialised functions or expression profiles exist, 
	we have curated and catalogued the entire human canonical core histone gene and protein complement.
	This work has led to multiple corrections to human histone annotations, 
	reflecting the flux of the human genome as it is updated and enriched in reference databases. 
	To avoid the obsolescence of gene catalogues but maintain readability of scientific literature, 
	we have developed a novel software framework that dynamically generates the up-to-date canonical histone repertoire 
	based on the most current genome annotations and inserts the information into a manuscript format. 
	This approach can readily be adapted to histones of other organisms and to other gene families.
  \end{abstract}

  \section{Introduction}
  %% Flow of introduction: histones are big part of the chromatin, the building unit.
  %% Then first division, split by type. When we know the type, we can explain
  %% nucleosome structure, which leads to core vs linker division. Then more history,
  %% we found variants. New division of classes, canonical vs variant. Explain the
  %% mess of variants, with all the names used for each division, and sometimes
  %% clashing nomenclature. Finalize with what this paper is about.

  Histones are among the most abundant proteins in eukaryotic cells and contribute up to half the mass of chromatin. 
  They define the structure and accessibility of the nucleosome as the fundamental repeating unit of genome organisation 
  around which the DNA is wrapped \addref{Luger1997 Nature}.
  In addition, the reactive side chains of histones are post-translationally modified 
  as a nexus for signalling and heritable epigenetics. 

  Relationships within the histone family have been described using a variety of terminologies
  reflecting their biochemical, functional and genomic distinctions (\tref{tab:histone-divisions})
  
  From a biochemical perspective, histones have been studied for many decades 
  as abundant proteins that are readily isolated using their highly basic chemical character. 
  Successive improvements in fractionation ultimately revealed 5 main histone types 
  with a nomenclature H1, H2A, H2B, H3 and H4 \citep{nomenclature}.
  An additional H1--related histone H5 is recognised in avian erythrocytes \citep{HFive-review}.

  The composition of arginine and lysine was an early distinction between the histones. 
  H1 has a low arginine/lysine ratio of \LinkerArgLysRatio{} and became known as ``lysine-rich'' 
  whereas the remaining 4 histones are ``arginine-rich'' with highest ratios in human isoforms of 
  \HTwoAArgLysRatio{},  \HTwoBArgLysRatio{}, \HThreeArgLysRatio{} and \HFourArgLysRatio{}
  for H2A, H2B, H3 and H4 respectively.

  The demonstration of the nucleosome as the fundamental repeating unit of chromatin \addref{Kornberg} 
  revealed that H2A, H2B, H3 and H4 associate as an octamer of two copies each within the
  nucleosome core particle, so they are referred to as ``core histones''. 
  H1 associates with the linker DNA between nucleosome core particles and is referred to as a ``linker histone''.

  Extracting histones and separating them by polyacrylamide gel electrophoresis (PAGE) 
  using the strongly anionic detergent sodium dodecyl sulphate (SDS) and neutral buffers 
  gives single bands for each histone type. 
  However, PAGE with non-ionic detergent Triton X--100 and uncharged urea as denaturants
  in acetic acid buffer (TAU or AUT) allows the separation of histones into multiple species 
  due to post-translational modifications and primary sequence variations within types \citep{PAGEND}.

  From a functional perspective, two classes of histone can be delineated 
  within each of the 4 core histone types (\tref{tab:typical-histone-differences}). 
  ``Canonical'' histones contribute to the bulk structure and generic function of chromatin. 
  These histones have highbut imperfect protein sequence identity.
  Their expression is significantly elevated during S~phase to provide chromatin packaging 
  for DNA duplicated during replication. This led to their description as ``replication dependent'', 
  although a supply of canonical histones is required throughout the cell cycle. 
  In contrast, a variety of ``variant'' histones such as H2A.Z, TH2B and H3.3 
  have less identity in alignments and lower aggregate abundance. 
  They play functionally distinct roles and are mostly expressed outside S~phase, 
  so are labelled as ``replication independent''. 
  Since nucleosomes containing variant histones are seen to be introduced in place of the canonical nucleosome, 
  the variants are also referred as ``replacement'' histones.
  
  From a genomic perspective the canonical histone genes are typically found in \todo{total histone clusters} clusters, 
  possibly for regulatory efficiency and to facilitate their sequence conservation. 
  These are somewhat confusingly referred to as ``allellic'' 
  and the histone gene product isoforms are often referred to as ``wild type'' although their equivalence has not been tested.
  The histone protein separations observed in TAU PAGE is sometimes acknowledged as ``sub-types'' 
  such as H2A.1, H2A.2 but this risks confusion with variants and belies their heterogeneity \addref{}. 
  Metazoan canonical histone genes are distinctive in lacking introns, 
  and giving rise to non-polyadenylated proten coding transcripts 
  where turnover is independently regulated via a highly conserved 3' stem-loop. 
  Most non-metazoan eukaryotes such as ciliate, fungi, and higher plants 
  all exhibit typical polyadenylated canonical histones transcripts \addref{}.

  %% I moderated the point about people mis-understanding H2A.2 etc
  
  Almost all variant histones are encode by single genes dispersed in the genome 
  with typical properties including introns, alternative splicing and polyadenylated transcripts, 
  so to contrast them with canonical histone genes they are referred to as ``non-allelic''
  

  %% I don't know enough to say there are no other non-polyadenylated protein transcripts so I washed this out
  %% I think the 3' stem loop is interesting but it is a distraction so let's leave it out
  %% Let's stay away from H1 because the danger is that we get sucked into including it too ...


  \begin{table*}
    \caption{General properties of canonical and variant histone proteins.}
    \label{tab:typical-histone-differences}
    \centering
    \begin{tabular}{l l l}
      \toprule
      \null                     & Canonical             & Variants \\
      \midrule
      Expression timing         & Replication dependent & Replication independent \\
      Sequence identity         & High                  & Low \\
      Functional relationships  & Isoforms              & Distinct specialised functions \\
      Transcript stabilisation  & Stem--loop in 3' UTR  & poly--A tail \\
      Gene distribution         & Clusters              & Scattered \\
      \bottomrule
    \end{tabular}
  \end{table*}


  \begin{table*}
    %% guess this could also be done in a description environment but...
    \caption{Typical sub-divisions of histones.}
    \label{tab:histone-divisions}
    \centering
    \begin{tabular}{F p{\dimexpr(\textwidth-\eqboxwidth{firstentry}-4\tabcolsep)}}
      \toprule
      Allelic variant & Same as canonical isoform. \\
      Canonical &   All of the histones listed in \tref{tab:histone-catalogue}.
                    Their main distinct features is a peak of expression during
                    S--phase and are usually regarded as without special
                    features. \\
      \addlinespace
      Heteromorphous & Same as variant. \\
      \addlinespace
      Homomorphous & Same as sub-type. \\
      \addlinespace
      Isoform   &   Often used incorrectly to mean a variant, it should be
                    used to mean an alternate form of the same histone type
                    either within the canonical or the same variant. \\
      \addlinespace
      Non-allelic variant & Same as variant. \\
      \addlinespace
      Replacement & Same as variant. \\
      \addlinespace
      Replication-dependent & Same as canonical. \\
      \addlinespace
      Replication-independent & Same as variant. \\
      \addlinespace
      Sub-type  &   H2A and H3 histones are sometimes further divided in the
                    subtypes H2A.1, H2A.2, H3.1, and H3.2. This division should
                    be avoided since it has no interpretation. Note that H3.3
                    is an histone variant, not a sub-type. \\
      \addlinespace
      Type      &   One of the 5 histone types, H1, H2A, H2B, H3, or H4.\\
      \addlinespace
      Variant   &   All of the histones listed in \tref{tab:variant-catalogue},
                    any histone that is not a canonical. Their main distinct
                    features is an expression independent of the cell cycle as
                    opposed to canonical histones. Sometimes, also incorrectly
                    named histone isoforms. \\
      \addlinespace
      Wild type &   Usually meaning the same as canonical or sub-type 1. This
                    naming is just wrong, since it suggests that variants are
                    not natural. \\
      \bottomrule
    \end{tabular}
  \end{table*}

  Despite the importance of histones to the organisation of chromatin and extensive interest
  in their role in epigenetics and regulation, the systematic curation and classification of human histone
  gene and protein sequences has not been revisited recently \citep{Marzluff02}.
  In the last decade the draft human genome has been finalised 
  and rich annotations continue to propagate into reference gene and protein sequence databases (\tref{tab:difference-from-Marzluff02}).

  %% This is primarily about canonical histones, not variants?

  In this paper we provide a comprehensive catalogue of the nomenclature and distribution 
  of canonical histone genes, proteins and pseudogenes. 
  Since curation and annotation are dynamic and evolving, we have implemented an accompanying resource 
  that can automatically regenerate this manuscript based on the most current data from the NCBI RefSeq database 
  in order to retain its value as a reference source in a scientifically accessible format.
  All automatically generated values are displayed with a light grey background.
  
  The seldom recognised diversity in canonical histones provides a foundation for investigating 
  their evolution, regulation and functional implications.



  \section{Histone genes}

  \subsection{Histone gene clusters}
	%% For simplicity, decided to list HIST4 explicitly for rather than allow it to come in by a variable
    The human canonical histone genes are located in clusters HIST1 to HIST4, 
	named in order of decreasing histone gene count (\tref{tab:histone-gene-count}).
	The major histone gene cluster HIST1 at chromosome~\HISTOneLocus{} and 
	contains \todo{percent genes in HIST1}\% of all histone genes 
	comprising \CodingGenesInHISTOne{}~functional core histone genes plus all canonical H1 histones. 
	HIST2 at chromosome~\HISTTwoLocus{} contains \CodingGenesInHISTTwo{}~coding genes. 
	HIST3 at \HISTThreeLocus{} and HIST4 at \HISTTFourLocus{} 
	contain only \CodingGenesInHISTThree{} and \CodingGenesInHISTFour{}~coding genes respectively. 
	Histone variants are dispersed across the genome outside the clusters, 
	so for example the three H3.3 encoding genes are located on chromosomes 1, 12 and 17 
	each far removed from other histone genes.

	HIST1 and HIST2 are both contiguous high density arrays of histone genes. 
	HIST1 spans \HISTOneSpan{} \todo{need space between number and unit}
	and is the second most gene dense region in the human genome at megabase scale \citep{MHC-III-analysis}.
	If \textit{HIST1H2AA} and \textit{HIST1H2BA} are excluded (see below), 
	the hemochromatosis protein encoding gene \textit{HFE} is the sole non-histone gene within HIST1. 
	There are no non-histone proteins within the HIST2 or HIST3 clusters. 

	The implications of this tight and exclusive clustering is not well understood \addref{Eirin-Lopez EvBiolCh8 2009}. 
	The dispersive organisation has been suggested to imply that expression is regulated by local elements \citep{close-regulators}, 
	and it is also possible that clustering influences histone gene sequence conservation. 
	HIST1 is located towards the distal end of the major histocompatibility complex (MHC) in the extended class~I region \citep{MHC-I-transcript, MHC-complete-sequencing-1999} 
	and it has been suggested that proximity to HIST1 suppresses recombination in the MHC Class~I region \citep{MHC-repressed-by-HIST}. 
	It has also been reported that HIST2 is prone to deletions and frequent rearrangements \citep{HISTTwo-prone-deletion-discovery, HISTTwo-prone-deletion-focus}. 

  \subsection{Histone gene nomenclature}
    Canonical histone gene names adhere to a Human Genome Organisation (HUGO) Gene Nomenclature Committee (HGNC) 
	endorsed system derived from the cluster number and position relative to other histones \citep{Marzluff02}.
	%% Should we try putting a url reference into the manuscript?
	%% http://www.genenames.org/genefamilies/histones

	The human canonical histone gene symbols are divided into 3 parts: 
	HIST cluster, histone type, and identifier letter
	for the order relative to other histone genes of the same type in the cluster (\fref{fig:nomenclature}A).
	For example, \textit{HIST1H2BD} is the fourth H2B gene in the HIST1 cluster.
	%% Changed example to match figure!
	The direction of the cluster order varies, so HIST1 is ordered from the centromeric side 
	whereas HIST2 and HIST3 are ordered from the telomeric side.
	%% Think this is correct from viewing the genome
	
    Two exceptions were originally applied to this simple gene naming rule \citep{Marzluff02}. 
	Firstly, the position identifier is omitted if there are no other histones of the same type in the cluster, 
	so \textit{HIST3H2A} is the sole H2A gene in HIST3. 
	Secondly, the human and mouse histone clusters are largely syntenic 
	so positional identifiers for missing orthologues were skipped to maintain the equivalence of gene symbols.
	This means there is no human \textit{HIST1H2AF} although \textit{Hist1h2af} exists in mouse.

	Several new histone genes have been uncovered since the original naming (\tref{tab:difference-from-Marzluff02}) 
	and this required additional nomenclature exceptions.
	For example, new H2A encoding genes identified in the human and mouse HIST2 cluster after \textit{HIST2H2AA} 
	obliged renaming of \textit{HIST2H2AA} to \textit{HIST2H2AA3} and the addition of the new gene as \textit{HIST2H2AA4}. 
	There are no human orthologues of mouse \textit{Hist2h2aa1} and \textit{Hist2h2aa2}.
	
	Furthermore, no distinction was originally made between pseudogenes and functional coding genes, 
	so \textit{HIST3H2AA} is a pseudogene whereas its neighbour \textit{HIST3H2AB} is a coding gene.
	Newly uncovered pseudogenes are now suffixed with PS and a number in order of discovery 
	instead of a positional identifier, such as \textit{HIST1H2APS5} 
	However, the ``original'' pseudogenes retain their symbols without PS. 
	This means that the absence of a PS suffix does not indicate a functional gene, 
	and that there is no positional information in the gene symbols of new pseudogenes.

	The flux in genome sequencing and increasing interest in histone genes 
	leaves open the possibility of further change. 
	For example, \textit{HIST1H2BA} is now recognised as testes-specific TH2 \addref{Horokoshi ACTA Cryst 2014?}, 
	has the most divergent sequence of canonical H2B proteins, 
	and is ~300 kbp upstream of the rest of the HIST1 cluster 
	separated by a number of unrelated non-histone genes. 
	\textit{HIST1H2BA} is adjacent to \textit{HIST1H2AA} 
	which encodes a similarly divergent protein compared the most common canonical H2A isoform.
	
	All recognised histone variant gene familes comprise only one or a few copies 
	dispersed across the genome outside clusters.
	The histone variant gene symbols have a separate 3 part nomenclature 
	consisting of labels for histone type, F for family, 
	and an identifier letter based on discovery order (\fref{fig:nomenclature}B).
	Where two or more isoforms exist the copies gain a numerical suffix, 
	although this is sometimes omitted for the first copy. 
	For example, \textit{H2AY} and \textit{H2AY2} encode the human macroH2As, 
	with \textit{H2AY} giving rise to alternative transcripts encoding macroH2A1.1 and macroH2A1.2 protein isoforms 
	whereas \textit{H2AY2} encodes macroH2A2.
	The variant histone gene symbols (\tref{tab:histone-variant-families}) also contain exceptions. 
	For example, the three gene copies encoding the H3.3 variant 
	take on distinguishing letters \textit{H3F3A}, \textit{H3F3B} and \textit{H3F3C} instead of numbers 
	and the H3 variant gene CENPA was named as the first identified centromeric protein \cite{CENPA-first-report} 
	before it was realised to be a histone protein \citep{CENPA-copurifies-histones, CENPA-sequence-analysis}. 
	Histone variant pseudogenes are suffixed with P to yield symbols such as \textit{H2AFXP} and \textit{H2AFZP5}.
	The increasing interest in histone variant function \addref{Maze Allis NRG 2014}
	but imprecise usage and conflicts between species, 
	has led to guidelines for improved consistency \addref{Talbert Epigen Chrom 2012}.

  \subsection{Histone gene sequences}

	The alignment of coding regions for functional canonical histone genes 
	demonstrates considerable sequence variation at teh DNA level (\fref{fig:histone-gene-variation}). 
	This relatively high sequence variation in human canonical histone genes 
	is consistent with observations across eukaryotes 
	that histone gene sequences are much less conserved than the proteins they encode 
	and that synonymous codon divergence exceeds non-synonymous variation \addref{Piontkivska MolBiolEvol 2000} \addref{Rooney MolBiolEvol 2002}. 
	Together this is taken as evidence that histone protein sequence conservation 
	results from strong selective pressure at teh protein level 
	and not from a homogenisation mechanism at the gene level based on clustering \addref{NeiRooney AnnRevGen 2005}.

    \begin{figure*}
      \centering
      \includegraphics[width=\textwidth]{nomenclature-schematic.pdf}
      \caption{Histone gene nomenclature. 
	  A. Canonical histone gene names encode their relative genomic order by cluster. 
	  B. Variant histone symbols are identified with F then a family identifier letter
	  and numbered by discovery, except 1 is omitted for the founder member. 
	  H3.3 genes conversely use numbers and letters for family and member.
	  C. Canonical seudogenes identified since 2002 are named by cluster with PS identifier numbered by discovery. 
	  Previously canonical pseudogenes were named like functional histone genes.}
      \label{fig:nomenclature}
	  \todo{Parts A and B need labels. Name text need to be italics!}
	  \todo{Change part B to use H2AFY2 as a more complete example of variant nomenclature}
	  \todo{Include part C using HIST1H2APS2 as example of PS pseudogene nomenclature}
    \end{figure*}

  \begin{table*}
    \caption{HGNC recognised histone variant family stem names, commonly used protein names 
	and names for improved consistency based on proposal by Talbert et al (2012). 
	HIST1H2BA is nominally not a variant but is included for completeness.}
    \label{tab:histone-variant-families}
    \centering
    \begin{tabular}{l l l}
      \toprule
      \Family & Protein name & Talbert et al name & Notes \\
      \midrule
	  H2AFB & H2A.Bbd & H2A.B & equivalent to H2A.L \\
	  H2AFJ & H2A.J & H2A.J & at HIST4 locus \\
	  H2AFV & H2A.Z-2 & H2A.Z.1 & - \\
	  H2AFX & H2AX & H2A.X & - \\
	  H2AFY & macroH2A/mH2A & macroH2A & - \\
	  H2AFZ & H2A.Z & H2A.Z.1 & - \\
	  H2BFM & ? & H2B.M & homologous to H2B.W, X-linked\\
	  H2BFS & - & - & pseudogene \\
	  H2BFWT & ? & H2B.W & testes specific, X-linked \\
	  H2BFX & - & - & pseudogene \\
	  HIST1H2BA & TH2B/TSH2B & H2B.1 & testes specific \\
	  H3F3 & H3.3 & H3.3 & euchromatin related \\
	  CENPA & CENP-A & ? & centromere-specific \\
	\bottomrule
    \end{tabular}
  \end{table*}

%% This is arguably the most important table. We should cite it as early as possible
    \begin{table*}
      \caption{
	  Count of human canonical histone coding and pseudogenes by histone cluster and type.
	  The right-hand value in each cell marked with $\psi$ is the number of pseudo genes.}
      \label{tab:histone-gene-count}
      \centering
      \begin{tabular}{l r@{ + }l<{$\psi$} r@{ + }l<{$\psi$} r@{ + }l<{$\psi$} r@{ + }l<{$\psi$} r@{ + }l<{$\psi$}}
        \toprule
        \null   & \multicolumn{2}{c}{H2A}  & \multicolumn{2}{c}{H2B}
                & \multicolumn{2}{c}{H3}   & \multicolumn{2}{c}{H4}
                & \multicolumn{2}{c}{Total} \\
        \midrule
        HIST1   & \HTwoACodingInHISTOne{}     & \HTwoAPseudoInHISTOne{}
                & \HTwoBCodingInHISTOne{}     & \HTwoBPseudoInHISTOne{}
                & \HThreeCodingInHISTOne{}    & \HThreePseudoInHISTOne{}
                & \HFourCodingInHISTOne{}     & \HFourPseudoInHISTOne{}
                & \CodingGenesInHISTOne{}     & \PseudoGenesInHISTOne{} \\
        HIST2   & \HTwoACodingInHISTTwo{}     & \HTwoAPseudoInHISTTwo{}
                & \HTwoBCodingInHISTTwo{}     & \HTwoBPseudoInHISTTwo{}
                & \HThreeCodingInHISTTwo{}    & \HThreePseudoInHISTTwo{}
                & \HFourCodingInHISTTwo{}     & \HFourPseudoInHISTTwo{}
                & \CodingGenesInHISTTwo{}     & \PseudoGenesInHISTTwo{} \\
        HIST3   & \HTwoACodingInHISTThree{}   & \HTwoAPseudoInHISTThree{}
                & \HTwoBCodingInHISTThree{}   & \HTwoBPseudoInHISTThree{}
                & \HThreeCodingInHISTThree{}  & \HThreePseudoInHISTThree{}
                & \HFourCodingInHISTThree{}   & \HFourPseudoInHISTThree{}
                & \CodingGenesInHISTThree{}   & \PseudoGenesInHISTThree{} \\
        HIST4   & \HTwoACodingInHISTFour{}    & \HTwoAPseudoInHISTFour{}
                & \HTwoBCodingInHISTFour{}    & \HTwoBPseudoInHISTFour{}
                & \HThreeCodingInHISTFour{}   & \HThreePseudoInHISTFour{}
                & \HFourCodingInHISTFour{}    & \HFourPseudoInHISTFour{}
                & \CodingGenesInHISTFour{}  & \PseudoGenesInHISTFour{} \\
        \addlinespace
        Total   & \HTwoACodingGenes{}       & \HTwoAPseudoGenes{}
                & \HTwoBCodingGenes{}       & \HTwoBPseudoGenes{}
                & \HThreeCodingGenes{}      & \HThreePseudoGenes{}
                & \HFourCodingGenes{}       & \HFourPseudoGenes{}
                & \TotalCodingGenes{}       & \TotalPseudoGenes{} \\
        \bottomrule
      \end{tabular}
    \end{table*}

    \begin{figure*}
      \centering
      \includegraphics[width=\textwidth]{seqlogo_canonical-coding-genes.pdf}
      \captionof{figure}{Sequence logos for canonical coding histone genes listed in table XX. 
	  Coding regions of A. H2A, B. H2B, C. H3 and D. H4.}
      \label{fig:histone-gene-variation}
    \end{figure*}

  \section{Histone transcripts}

  Canonical core histone transcripts are among the only protein coding messages
  that do not undergo polyadenylation,
  instead carrying a unique stem-loop structure (\tref{tab:typical-histone-differences}).
  They do have a 7--methyl--guanosine cap at the 5' end \citep{MarzluffNatRevGen2008}.

  Canonical core histone gene transcription is regulated
  by cell-cycle dependent phosphorylation of the histone-specific 
  Nuclear Protein of the Ataxia-Telangiectasia locus (NPAT) coactivator
  and interaction with the accessory protein 
  FADD-Like interleukin-1$\beta$-converting enzyme/caspase-8-ASsociated Huge protein (FLASH),
  resulting in assembly of histone locus bodies
  coordinating factors responsible for transcription and processing
  \citep{MarzluffNatRevGen2008,RattrayMueller2012,Hoefig2014}.
  Variability is observed in canonical core histone isoform gene transcription,
  both by analysis of non-polyadenylated transcripts \citep{YangGenomeBiol2011}
  and RNA polymerase II promoter occupancy \citep{Ederveen2011}.

  Overall there is estimated to be a 35 fold increase in mammalian
  canonical histone transcripts during S phase,
  principally as a result of a 10 fold increase in mRNA stabilisation
  via stem-loop dependent mechanisms,
  and a 3--5 fold up-regulation in canonical core histone gene transcription \citep{HarrisMCB1991}.

  This post-translation regulation of canonical core histone transcripts
  is achieved by a stem-loop encoded in the mRNA 3' untranslated region
  that is recognised by specialised spliceosome-related RNA
  processing and stabilisation complexes \citep{stem-loop-structure}.
  The modal location of annotated stem-loops is \StemLoopStart{} to \StemLoopEnd{} bp after the stop codon
  for the set of all canonical histone genes.
  The sequence logo of all aligned stem-loops
  confirms these are highly conserved in sequence and structure (\fref{fig:stem-loop-seqlogo}).

  The RNA stem-loop structure is bound by the Stem-Loop Binding Protein (SLBP)
  which is up-regulated 10--20 fold during S~phase to stabilise
  histone mRNAs \citep{SLBP-regulation}.
  Immediately downstream of the stem-loop a purine-rich Histone Downstream Element (HDE)
  interacts with U7 snRNA to direct efficient 3' end processing.
  Although this is not an annotated feature of histone genes,
  alignment of the canonical HDE \citep{HDE-sequence} to all canonical histone genes
  shows the modal location of the HDE is 
  \HDEsDistanceFromStemLoop{} bp downstream of the stem-loop.
  The sequence logo calculated from all aligned HDEs 
  confirms that this feature is also highly conserved (\fref{fig:HDE-seqlogo}).

  \begin{figure}
    \centering
    \subtop[Stem-loop]{
      \includegraphics{seqlogo_stem_loops.pdf}
      \label{fig:stem-loop-seqlogo}
    }
    \hfill
    \subtop[HDE]{
      \includegraphics{seqlogo_HDEs.pdf}
      \label{fig:HDE-seqlogo}
    }
    \caption{%
      Sequence logos for A) annotated stem-loops, and 
      B) Histone Downstream Elements (HDEs) identified by homology
      for all canonical core histone gene 3' untranslated regions (UTRs).
    }
  \end{figure}

  Stabilisation and processing extend the half life of canonical core histone \mbox{mRNAs}
  and contribute to increased histone translation efficiency,
  enabling rapid production of histones to package the newly duplicated genome.

  Although the vast bulk of canonical core histone transcripts 
  appear to be regulated by this mechanism, 
  1-5\% of transcripts are found to be 3' polyadenylated \citep{YangGenomeBiol2011}
  and some genes have annotations indicating two transcripts
  differing in whether they have stem-loop or polyadenylation signals (\tref{tab:curation-anomalies}).

  Core histone variant transcripts lack a 3' stem-loop and are
  polyadenylated in the same way as almost all other protein coding genes.
  The exception is H2AX, which has alternatively processed transcripts
  exhibiting the stem-loop characteristic of canonical core histones
  or a poly(A) tail found on variant core histones \citep{HTwoAX-transcripts,our-H2AX-review}.

  \section{Histone proteins}
  Most depictions of chromatin imply that
  canonical core histone protein types behave as a single protein,
  often referred to as ``wild type''.
  The assumption is based on the relatively high identity of histone sequences
  between isoforms and species,
  and the focus of interest on specific functional roles of core histone variants.

  However, the very strong selection pressure on
  amino acid sequences encoded by canonical core histone genes \citep{NeiRooney2005},
  the roles of specific amino acid differences in observed proteins \citep{MazeAllis2014},
  and the consequences of small variations in proteins on the structure of nucleosomes \citep{KurumizakaCOSB2013}
  all suggest that the minor differences in the encoded canonical core histone protein isoforms
  could have functional implications.

  Encoded isoform variation in canonical histone types H2A and H2B is most pronounced,
  with \HTwoAUniqueProteins{}~H2A and \HTwoBUniqueProteins{}~H2B distinct polypeptides
  encoded by these histone type genes.
  In contrast, only \HThreeUniqueProteins{}~H3 and \HFourUniqueProteins{}~H4
  distinct polypeptides are encoded by similar numbers of genes (\tref{tab:histone-gene-count}).
  Although most variation is a result of amino acid substitutions,
  some H2A and H2B isoforms also show C-terminal length differences.

  The nomenclature for histone protein isoforms follows directly
  from the gene names \citep{Marzluff02},
  and supersedes the earlier nomenclature with forward slash which used different isoform letters
  \citep{AlbigGenomics1997,AlbigHumangen1997}.
  The encoded proteins described below are products of genes and transcripts listed
  in table \tref{tab:histone-catalogue}.
  Human canonical core histone polypeptides have typically been numbered
  with omission of the N-terminal methionine
  since this is likely to be removed because most sequences have
  a small hydrophilic amino acid as the second residue \citep{XiaoPeiBiochem2010}.
  This convention predates the Human Genome Variation Society (HGVS) recommendation
  to include the initial methionine as residue 1.
  We have omitted the N-terminal methionine on figures, alignments,
  and amino-acid numbering for consistency.

%% To include methionine just uncomment 3 lines in the routine that reads the sequence files in MyLib.pm

  \subsection{Canonical H2A isoforms}
    Canonical H2A genes encode \HTwoAUniqueProteins{}~different H2A
    protein isoforms with pairwise identity down to
    \FPround{\result}{\HTwoAPID}{0} \SI{\result}{\percent} (\tref{tab:H2A-consensus}).
    These are separable by TAU PAGE into two bands identified as H2A.1 and H2A.2.
    TAU PAGE distinguishes leucine or methionine at residue 51 \citep{FranklinZweidler1977,Zweidler1977},
    implying there are up to 8 different protein sequences in H2A.1 and 3 in H2A.2.
    There is no concordance between H2A.1 and H2A.2
    and the location of their isoform-encoding genes in the histone gene clusters.
    The HIST2 cluster contains genes encoding isoforms with both Leu51 and Met51,
    while HIST1, HIST2, and HIST3 clusters all contain genes encoding H2A isoforms with Leu51.
    No functional distinction between H2A.1 and H2A.2 has yet been reported.

    Excluding \textit{HIST1H2AA} discussed above,
    the sites of difference in two or more isoforms of canonical H2A are
    serine or threonine at residue 16,
    alanine or serine at residue 40,
    lysine or arginine or glycine at residue 99,
    and the C-terminal residues from residue 124 onwards (\tref{tab:H2A-consensus}).
    All these sites have implications for post-translational modifications.

    \begin{table}
      \caption{%
        Canonical H2A encoded protein isoforms.
        Upper panel shows isoform variations relative to most common isoform
        using HGVS recommended nomenclature \citep{mutnomenclature2003}.
        Lower panel shows sequence logo of all isoforms aligned
        with invariant residues in grey.
      }
      \label{tab:H2A-consensus}
      \input{results/table-H2A-proteins-align}
      \includegraphics[width=\textwidth]{seqlogo_H2A_proteins.pdf}
    \end{table}

  \subsection{Canonical H2B isoforms}
    Canonical H2B has more isoforms than the other histone types (\tref{tab:H2B-consensus})
    with \HTwoBUniqueProteins{}~unique proteins
    diverging to \FPround{\result}{\HTwoBPID}{0}\result\% identity.
    Nevertheless, all isoforms appear to migrate together in TAU PAGE.

    There is significant variability between isoforms in the N-terminal region
    and this is one of the most variable sites between canonical
    core histones of different species.
    Human H2B has a unique and distinctive N-terminal proline-acidic-proline motif (PEP/PDP)
    followed by a very variable residue with either alanine, serine,
    threonine, or valine (\tref{tab:H2B-consensus}).

    Excluding \textit{HIST1H2BA} discussed above,
    the remaining isoform differences in two or more isoform are mainly the chemically
    similar valine or isoleucine at residue 39,
    and serine to glycine and alanine at residues 75 and 124 respectively,
    both of which have post-translational modification implications.
    A number of H2B genes are annotated to have multiple transcripts,
    although in most cases these transcripts encode identical protein isoforms.
    Variation in H2B transcript isoform levels between multiple cancer cell
    lines has been observed \citep{Molden2015},
    although the functional implications are unknown.

    \begin{table}
      \caption{%
        Canonical H2B encoded protein isoforms.
        Upper panel shows isoform variations relative to most common isoform
        using HGVS recommended nomenclature \citep{mutnomenclature2016}.
        For clarity, isoforms encoded by multiple transcripts of a single gene
        are distingushed by a numerical suffix (see \tref{tab:histone-catalogue}).
        Lower panel shows sequence logo of all isoforms aligned
        with invariant residues in grey.
      }
      \label{tab:H2B-consensus}
      \input{results/table-H2B-proteins-align}
      \includegraphics[width=\textwidth]{seqlogo_H2B_proteins.pdf}
    \end{table}

  \subsection{Canonical H3 isoforms}
    Canonical H3 genes encode only \HThreeUniqueProteins{}~different
    protein isoforms (\tref{tab:H3-consensus}).
    The majority of H3 genes are in the HIST1 cluster and encode a
    single polypeptide sequence \citep{Ederveen2011}
    whereas the three canonical H3 genes in HIST2 encode a distinct isoform
    with the interesting difference of serine instead of cysteine at residue 96
    that is separable by TAU PAGE \citep{FranklinZweidler1977}.
    By apparent coincidence this means HIST1-encoded canonical H3 isoforms are
    identified as H3.1 and the HIST2-encoded copies are H3.2.
    As discussed above, \textit{HIST3H3} with four amino acid differences
    appears to be the largely testes specific variant H3T
    that has also been assigned as H3.4 in variant nomenclature \citep{Talbert2012}
    even though it would be predicted to migrate with H3.1 in TAU PAGE \citep{FranklinZweidler1977}.

    \begin{table}
      \caption{%
        Canonical H3 encoded protein isoforms.
        Upper panel shows isoform variations relative to most common isoform
        using HGVS recommended nomenclature \citep{mutnomenclature2003}.
        Lower panel shows sequence logo of all isoforms aligned
        with invariant residues in grey.
      }
      \label{tab:H3-consensus}
      \input{results/table-H3-proteins-align}
      \includegraphics[width=\textwidth]{seqlogo_H3_proteins.pdf}
    \end{table}

  \subsection{Canonical H4 isoforms}
    H4 is the most homogeneous of all canonical core histones,
    with all but one of the \HFourCodingGenes{}~genes encoding
    identical proteins (\tref{tab:H4-consensus}).
    These genes are located across HIST1, HIST2, and the isolated
    \textit{HIST4H4} as the only gene in HIST4.

    The divergent \textit{HIST1H4G} gene in the middle of the HIST1 cluster
    encodes an isoform with 15 amino acid differences and a deletion of the C-terminal 5 residues.
    This gene is annotated as being transcribed and merits further investigation.

    \begin{table}
      \caption{%
        Canonical H4 encoded protein isoforms.
        Upper panel shows isoform variations relative to most common isoform
        using HGVS recommended nomenclature \citep{mutnomenclature2003}.
        Lower panel shows sequence logo of all isoforms aligned
        with invariant residues in grey.
      }
      \label{tab:H4-consensus}
      \input{results/table-H4-proteins-align}
      \includegraphics[width=\textwidth]{seqlogo_H4_proteins.pdf}
    \end{table}


    %% TODO
%  \section{SNP's}

  \section{Canonical isoforms vs.~variants}
  %% Characteristics of a variant are the fact they seem to have special functions. Emphasis
  %% that some canonical may have special functions which no one bothered to check
  %% variants with extra special nomenclature are outside the histone clusters and have poly-A tails (sometimes)
  %% Complain about the lack of information on the different isoforms

  %% Should we be pointing the finger to bad examples? That's really bad form.
  %% We could show examples where H2AX is called histone isoform, or even
  %% canonical histones are named variants, but really, what's the pointing of
  %% pointing the fingers to those who are doing it wrong other than shame them?

  The terms canonical, variant, and isoform histones are often misused, although
  the correct meaning can usually be infered from the context. This
  confusion is likely to stem from not knowing the existence of the two groups,
  rather than confusion about where they belong.

  There are canonical and variant histones, each group listed in \tref{tab:histone-catalogue}
  and~\ref{tab:variant-catalogue} respectively. Within these two groups,
  there are also isoforms. For example, each of the \HTwoAUniqueProteins{} unique H2A proteins
  is an isoform of the canonical H2A, while the multiple H2AFY are isoforms of
  the H2AFY protein.

  We have summarized the distinction between canonical and variants in
  \tref{tab:typical-histone-differences} which is then reflected on their
  gene symbols (\fref{fig:nomenclature}). While this separation appears like a
  good organization, findings of new genes and studies start to stress those
  definitions to the point where their distinction is unclear. And since
  gene symbols should be stable, such distinction is actually harmful. We will
  either end up with changing or meaningless symbols.

  %% "Attack" each of the identifying characteristics

  %% The paper describing H2AFJ, with the RT-PCR has a plotthat shows H2AFJ
  %% levels stable while HIST2H2AA2 increases during S phase. But shouldn't
  %% canonical histones increase 35 times instead of only 2.3? On the other
  %% hand, the plot does not have legend for either the X or Y axis, I can
  %% probably be interpretting it wrong.

  The main distinction between the two groups is related to its dependence on
  cell cycle. Canonical histones are mainly expressed during S~phase due to
  its mRNA 3' processing via the stem-loop mechanism. This has been used
  as the single deciding characteristic used to name H2AFJ as an histone variant,
  despite no knowledge of specialized function\citep{h2afj-report}. In the study,
  it is shown that H2AFJ has a poly-A tail instead and not a stem-loop, and that
  its expression levels do not increase during S--phase.

  However, it is now also argueed that some canonical H2B histones genes, encode
  two transcripts, with a stem-loop and with a poly(A) tail. This suggest
  that not all of the canonical histones are completely replication dependent.
  An equivalent system is used by the variant H2AX which is usually regarded
  as replication independent.

  %% TODO the sanity check list at the end of the paper, should be made into a
  %%      a table, or something else that can be referenced, so we can use it
  %%      in the text

  In addition, the variant vs.~canonical separation suggests that there's no variance inside
  the canonicals which is false. There's several unique protein sequences for each
  `canonical' protein. And even the sometimes named H2A.2 and H3.2 variants are actually
  canonical histone proteins, with no more sequence variance than the other canonical
  isoforms.

  Finally, the word ``canonical'' is suggests precedence, the original histones
  from which variants diverged. We have suggested, as others did, that this
  is not necessarily true, and that it may have been the H2AX ``variant'' the
  original from which the ``canonical'' H2A evolved\citep{our-H2AX-review, henikoff2010-variants-evolution}.


  Variants are the most likely with a standard gene, are considered independent from each other,
  spread around the genome and its expression is mostly independent of cell cycle \todo{explain
  why mostly? Do it next paragraph when complaining the naming}. When people say H2AX that's
  a variant. Note, H2A.2 and H3.2 are not considered variants.


  However, the term ``variant histones'' implies that canonical histones are
  homogeneous and masks their heterogenity since for sure, if another histone
  of the same type is different, it ought to be a variant. This is not
  true. Each variant, such as H2A.Z, H2Abbd, and H3.3,
  is encoded by a gene, with its own sequence and name,
  but canonicals are encoded by multiple genes with less varying sequence.

  It is often said in software development, that no documentation is better
  than incorrect documentation. We stand by this opinion and arguee that
  a nomenclature that is sometimes incorrect does more wrong than right.
  Incorrect documentation of software
  is usually the result of evolving code whose documentation is left behind,
  not of documentation incorrectly written at the start.
  This is the same with histone gene nomenclature. It correctly delivered
  the, now outdated, information about them. However, as the knowledge of
  these genes increased,
  their symbols simply become misleading and a source of wrong assumptions.

  A list of all the changes since \citet{Marzluff02} is
  displayed on \tref{tab:difference-from-Marzluff02}.


  With this article being closer to a software release than a paper publication,
  patches are welcome.


  when the gene
  symbol is given there is no doubt ever.

  \section{Reproducible research}

  A considerable number of builds and updates of the human genome sequence have been made
  since the last major survey of human canonical histone genes in 2002 \citep{Marzluff02}.
  This has resulted in
  \AddedSinceReference{} canonical core histone genes added,
  \RemovedSinceReference{} removed,
  and \SequencesUpdatedSinceReference{} sequences
  being updated for the current RefSeq release (\tref{tab:difference-from-Marzluff02}).
  These changes reflect the ongoing nature and challenges of
  genome curation and annotation in reference databases \citep{BorkKoonin1998}.

  It is inevitable that sequences and annotations will continue to be revised
  based on continuous curation and improved experimental insights.
  This in turn prompts reevaluation of assumptions about their biological contributions,
  illustrated by the recognition that some genes
  annotated as canonical core histone isoforms are testes-specific variants \citep{Talbert2012},
  and by new observations of cell type-specific expression \citep{Molden2015}.

  It is important for communities of researchers to contribute
  to ongoing formal curation of genomics resources
  by feeding back information to database maintainers \citep{SteinNRG2001},
  although many are unaware of this opportunity \citep{HollidaySPR2015}.
  In the course of this work we have suggested a significant number of proposals
  for improvements to RefSeq curators.
  These were identified automatically by scripted tests for consistency and uniformity
  between the database annotations and expected properties of histones.
  \tref{tab:curation-anomalies} lists the current apparent anomalies from these tests
  and is the basis for ongoing discussions with curators.

  Dynamic data can be presented in specialist online database resources
  such as the Histone Database \citep{HistoneDB2016} and HIstome \citep{HIstome2012}.
  These database interfaces provide comprehensive access to data
  but have limited curation and are difficult to cite.
  They can also cease to be updated or become inaccessible.
  In contrast, manuscripts are the established method of scientific communication
  because they provide descriptive context and accessible formatting for readers,
  can be peer reviewed, and have well-established methods for citation.
  There are established mechanisms for permanent archival of published manuscripts.
  However, static catalogues in manuscripts
  such as earlier surveys of histone genes \citep{AlbigHumangen1997,Marzluff02}
  inevitably become supersed by improvements in source data and cannot be updated to remove errors.

  A self-updating manuscript bridges the features of dynamic database and static manuscript
  presentation styles by providing convenient access to the most current information.
  It encourages communities of researchers to directly feed into the formal curation process
  and enables rapid leveraging of the most current data for relatively stable gene families.

  Implementating a dynamic manuscript is also an example of ``reproducible research''
  \citep{reproducible-research-bioinformatics,reproducible-research-law}
  that can address the topical challenge of irreproducibility in biological data
  and interpretation \citep{ErrorProne2012,OpenPrograms2012}.

  This manuscript is generated directly from sequences and annotations
  in the core NCBI RefSeq resource \citep{OLearyRefseq2016}.
  The processing system for the manuscript does not cache intermediate information,
  so all changes contributed by the community of histone researchers
  and curated by professional database maintainers
  are directly and automatically reflected at each manuscript refresh.
  All scripts including instructions for automatic builds
  are transparently available in a public repository.
  Core processing is based on a BioPerl program contributed to the Bio-EUtilities distribution
  and publicly available since 2013.
  Dependencies on raw data sources, alignment algorithms, or display output libraries
  can also be upgraded since the processing is automatic and script-based.

  One potential challenge for researchers is the referencing of dynamic data within such a manuscript.
  This can be simply addressed by users citing the publication as a traditional static manuscript
  then stating the build date in Materials and Methods, as they would for a database.

  Although the dynamic data will remain current,
  major revisions in understanding of a gene family or accumulation of small insights
  can in time render the explanatory static manuscript text obsolete.
  In this case the manuscript can simply be refreshed and republished independently
  in the same way as a traditional manuscript would be.
  The underlying scripts generating the dynamic data do not need to be rewritten,
  although new functions can be added to reflect new insights.

  The core scripts underlying this manuscript have been written to facilitate generating
  equivalent catalogues for other organisms via simple build options.
  We have successfully trialled this for \textit{Mus musculus} (not shown)
  to demonstrate that an equivalent dynamic manuscript for other histone gene sets
  only requires drafting appropriate surrounding static explanatory text.
  The approach is also applicable to larger and more diverse gene families such as
  our previous cataloguing of the Snf2 family \citep{andrew-snf2-catalogue}.



  \section{Acknowledgements}
    We are grateful to all the nice people at \#perl and \#bioperl channels on
    \url{irc.freenode.net}, especially Altreus, mst, LeoNerd, and pyrimidine for
    their help with the most technical programming details.

  \bibliographystyle{plainnat}
  \bibliography{references}

  %% Supplementary tables and figures
  \onecolumn

  \newpage
  \captionof{table}{Differences between the current reference (as per RefSeq)
                    for human histone genes and the last published reference \citep{Marzluff02}.}
  \label{tab:difference-from-Marzluff02}
  \input{results/table-reference_comparison}

  %% this is a very long table (ctabular) that does not play nice with two columns
  %% so we place it at the end. See why it was a problem at
  %% http://tex.stackexchange.com/questions/101791/float-a-wide-ctabular-in-two-columns
  %% We could split the table in two parts in which case we could use begin{table*}
  \newpage
  \captionof{table}{Catalogue of all human canonical core histone genes, with
                    respective transcripts, protein, and IDs}
  \label{tab:histone-catalogue}
  \input{results/table-histone_catalogue}

  \newpage
  \captionof{table}{Catalogue of all human variant core histone genes, with
                    respective transcripts, protein, and IDs}
  \label{tab:variant-catalogue}
  \input{results/table-variant_catalogue}

  \newpage
  \input{results/histone_insanities.tex}

\end{document}
